\documentclass[11pt]{article}
\usepackage[margin=1.1in]{geometry}
\usepackage{amsmath,amssymb,amsthm,hyperref,booktabs}
\newtheorem{theorem}{Theorem}
\title{The $10\times42$ Chomp bar has three winning moves}
\author{Jaideep Sai Padhi\\ \small Purdue University \quad \texttt{jpadhi@purdue.edu}}
\date{August 2026}
\begin{document}\maketitle
\noindent\textbf{Keywords:} Chomp; combinatorial games; winning-move multiplicity; retrograde analysis.\\
\textbf{MSC 2020:} 91A46 (primary); 05-04, 68R05 (secondary).
\begin{abstract}
In Chomp on an $a\times b$ chocolate bar, Gale's strategy-stealing argument shows the first
player wins, but exhibits no winning move; famously, winning moves are known only by exhaustive
computation, and their number varies. Ken Thompson found in the 1970s that the $8\times10$ bar
has two winning moves; Ekhad and Zeilberger (2018), tabulating all bars up to $14\times14$,
found further doubles and challenged: find a bar with at least \emph{three} winning moves. We
answer the challenge: \textbf{the $10\times42$ bar has exactly three winning bites, at
$(5,36)$, $(7,30)$ and $(8,26)$.} The computation is an exhaustive retrograde solve of the
full ideal lattice, verified by two independent implementations and against the published
tables. We also extend the winning-move-multiplicity census far beyond the 2018 range
(eighteen bars with exactly two winning moves in the searched region), prove a small
structural theorem --- the winning bites of any bar form a strict staircase --- and report
striking quasi-periodic structure in where multiple winning moves occur, including
bounded-error Beatty-type laws with rotation numbers in $\mathbb{Q}(\sqrt2)$ for the
diagonal families of three-row Chomp.
\end{abstract}

\section{The challenge}
Chomp is played on an $a\times b$ grid of cells (rows indexed downward); a move (a
\emph{bite}) at cell $(i,j)$ removes every cell $(r,c)$ with $r\ge i$ and $c\ge j$; the
player forced to take the poisoned cell $(1,1)$ loses. Positions are exactly the Young-diagram
ideals $\lambda=(\lambda_1\ge\cdots\ge\lambda_a)$. Gale's strategy-stealing argument
\cite{Gale} shows every rectangular bar (other than $1\times1$) is a first-player win, but
famously certifies no winning move. Computation is the only known access, and already Gale asked whether the winning first
move is always \emph{unique}. The answer is no --- Ken Thompson discovered that the
$8\times10$ bar has two winning moves (see \cite{BHMS}) --- but the uniqueness phenomenon is
real in low rows: Sheiner \cite{Sheiner} has recently proved that every $3\times n$ bar has
exactly one winning move, settling the three-row case of Gale's question, and Garg
\cite{Garg} has tabulated the $4\times n$ P-positions to $n\le500$. Ekhad and Zeilberger
\cite{EZ}, computing all bars up to $14\times14$, found a handful of doubles beyond
Thompson's ($6\times13$, $9\times10$, $10\times14$, $12\times13$) and posed as their Second
Computational Challenge: \emph{find $a,b$ such that the $a\times b$ bar has at least three
winning moves.} This note answers the challenge and maps the anti-uniqueness landscape.

\begin{center}
\fbox{\parbox{0.86\textwidth}{\textbf{Answer.} The $10\times42$ bar has exactly three winning
bites: at $(5,36)$, $(7,30)$, and $(8,26)$.}}
\end{center}

\section{A staircase theorem}
\begin{theorem}\label{thm:stair}
If $(i,j)$ and $(i',j')$ are both winning bites of the same bar and $i<i'$, then $j>j'$.
Moreover no winning bite of an $a\times b$ bar with $a\ge2$, $ab\ge2$ lies in row $1$.
\end{theorem}
\begin{proof}
The bite $(i,j)$ on the full bar leaves the two-block position $B(i,j)$ with rows
$1,\dots,i-1$ of length $b$ and rows $i,\dots,a$ of length $j-1$. If $i<i'$ and $j\le j'$
then biting $B(i',j')$ at $(i,j)$ yields exactly $B(i,j)$: two winning bites leave
P-positions, and no P-position is one move from another. For the second claim, a row-1 bite
leaves a full $a\times(j-1)$ rectangle, an N-position by Gale (or, for $j=2$, a column, which
is N for $a\ge2$).
\end{proof}
The three bites of $10\times42$ --- rows $5<7<8$, columns $36>30>26$ --- form the first
three-step staircase ever observed. Empirically the staircase is usually \emph{tight} --- in sixteen of the eighteen known
doubles, and in the triple, consecutive winning bites are $1$ or $2$ rows apart --- but not
always: the two side-$16$ doubles split by six rows (Table~\ref{tab:bites}).

\section{Algorithm}
One retrograde solve of the ideal lattice of an $A\times B$ box determines the winning-move
count of \emph{every} bar $a\times b$ with $a\le A$, $b\le B$, since each such bar is a
position in the box. We rank positions by the combinatorial number system on
$d_k=\lambda_k+A-k$ (strictly decreasing), i.e. $\mathrm{rank}(\lambda)=\sum_k
\binom{d_k}{A-k+1}$: this bijection onto $\{0,\dots,\binom{A+B}{A}-1\}$ is monotone under
componentwise order, so a single increasing-rank sweep is a valid dynamic-programming order.
The children of $\lambda$ under a bite in row $i$ at threshold $t$ (all rows $\ge i$ capped
at $t$) have ranks computable in $O(1)$ amortized via the hockey-stick identity
$\sum_{m=m_1}^{m_2}\binom{t-1+m}{t-1}=\binom{t+m_2}{t}-\binom{t+m_1-1}{t}$ applied to the
capped rows. The production engine stores one bit per position (P/N), detects N-positions
with early exit at the first winning move, recovers exact counts for the bars in a final
pass, and parallelizes by area layers (all positions of area $s$ depend only on smaller
areas). On $13$ cores it solves the $\binom{55}{10}\approx2.9\times10^{10}$-state
$10\times45$ box --- the box that found the triple --- in under half an hour.

\section{Validation}
(i) All engines agree with brute-force game search on six small boxes, both per-bar
winning-move counts and the full P-position census. (ii) The published record is reproduced
exactly: Thompson's $8\times10$ has two winning moves, at $(4,9)$ and $(5,6)$, and the
Ekhad--Zeilberger doubles $6\times13$, $10\times14$ check. (iii) The triple was found in the
$10\times45$ box, re-derived from the (different) $10\times42$ box, and counter-signed by an
independent implementation with a different iteration order and no early exit. All code and
logs: \url{https://github.com/jaideepsaipadhi/chomp-three-winning-moves}.

\section{The census}\label{sec:census}
In the searched region --- all $a\le b\le19$; $4\times b$, $b\le300$; $5\times b$, $b\le150$;
$6,7\times b$, $b\le110$; $8\times b$, $b\le60$; $9\times b$, $b\le55$; $10\times b$,
$b\le45$ --- there are exactly \textbf{eighteen} bars with exactly two winning moves and
\textbf{one} with three:
\begin{center}
$6\times13,\ 6\times93,\ 7\times29,\ 7\times30,\ 7\times57,\ 8\times10,\ 8\times22,\
8\times23,\ 9\times10,$\\ $9\times26,\ 10\times14,\ 10\times29,\ 10\times33,\ 10\times35,\
11\times18,\ 12\times13,\ 13\times16,\ 14\times16$
\end{center}
(and transposes); every other bar in the region, including every square through $19\times19$,
has a unique winning move; in particular $4\times b$ and $5\times b$ bars are unique
throughout their ranges, consistent with (and, for four rows, implicit in the data of)
Garg's tabulation \cite{Garg}.
Three regularities stand out. \emph{Adjacent pairs:} doubles arrive at consecutive widths
($7\times29/30$, $8\times22/23$) or consecutive heights ($13\times16$, $14\times16$).
\emph{Hot and barren lines:} no bar with a side of $15$ has a multiple move, while side $16$
carries two doubles; row $10$ carries six doubles \emph{and} the triple. \emph{Long gaps:}
row $6$ is silent from $14$ to $92$, then hits at $93$.

\begin{table}[h]\centering
\caption{Winning bites of every known multi-move bar (staircase form guaranteed by
Theorem~\ref{thm:stair}); the crossing families $(p,q)=(i{-}1,a{-}i{+}1)$ can be read off
the bite rows.}\label{tab:bites}
\begin{tabular}{llll}
\toprule
bar & bites & bar & bites\\
\midrule
$6\times13$ & $(4,12),(5,9)$ & $10\times14$ & $(5,13),(7,9)$\\
$6\times93$ & $(2,52),(3,39)$ & $10\times29$ & $(6,16),(8,12)$\\
$7\times29$ & $(6,25),(7,20)$ & $10\times33$ & $(7,23),(9,18)$\\
$7\times30$ & $(4,21),(5,20)$ & $10\times35$ & $(7,24),(8,21)$\\
$7\times57$ & $(3,42),(4,26)$ & $11\times18$ & $(9,17),(10,16)$\\
$8\times10$ & $(4,9),(5,6)$ & $12\times13$ & $(10,11),(11,9)$\\
$8\times22$ & $(2,9),(3,7)$ & $13\times16$ & $(5,6),(11,5)$\\
$8\times23$ & $(7,18),(8,16)$ & $14\times16$ & $(3,12),(9,10)$\\
$9\times10$ & $(7,8),(9,5)$ & $\mathbf{10\times42}$ & $\mathbf{(5,36),(7,30),(8,26)}$\\
$9\times26$ & $(2,11),(4,10)$ & & \\
\bottomrule
\end{tabular}
\end{table}

\section{The structure of multiplicity}
The count decomposes: biting the $a\times b$ bar at $(i,j)$ leaves the two-block position
with $p=i-1$ full rows and $q=a-p$ rows of length $t=j-1$, so the number of winning moves of
the bar equals the number of families $(p,q)$, $p+q=a$, whose \emph{P-track}
$\{(w,t):\ (w^p,t^q)\text{ is a P-position}\}$ passes through width $w=b$. Empirically each
track is sparse and quasi-linear --- essentially one P-cell per width --- so doubles are
crossings of two tracks and the triple is a threefold crossing, constrained by
Theorem~\ref{thm:stair} to a staircase. The two $7$-row doubles, for instance, are the
crossings $(5,2)$--$(6,1)$ at $b=29$ and $(3,4)$--$(4,3)$ at $b=30$.

For three-row Chomp the tracks obey laws of unexpected exactness. Writing the $n$-th P-cell
of the family $(2,1)$ (positions $(w,w,t)$) as $(w_n,t_n)$, we find, over $247$ cells up to
$w=600$,
\[ w_n=(1+\sqrt2)\,n+O(1),\qquad t_n=\tfrac{2+\sqrt2}{2}\,n+O(1), \]
with error confined to a window of width $<3$ and \emph{no drift}; the family $(1,2)$
(positions $(w,t,t)$) hits essentially every $t$, with $w(t)=\tfrac{2+\sqrt2}{2}\,t+O(1)$,
window $<2.4$ over $351$ cells. The slopes' continued fractions match $\sqrt2=[1;\overline2]$
and $\tfrac{2+\sqrt2}{2}$ through all computed partial quotients. We state these as
conjectures: the diagonal families of three-row Chomp are quasi-periodic with rotation
numbers in $\mathbb{Q}(\sqrt2)$. The exact three-row machinery of Brouwer et
al.\ \cite{BHMS}, in which Sheiner's uniqueness theorem \cite{Sheiner} is proved, is the
natural arena in which to attack them. This is finer than, and consistent with, the renormalization
picture of Friedman and Landsberg \cite{FL}, who analyzed the three-row P-region
probabilistically; Zeilberger's assessment of three-row Chomp as chaotic \cite{Zchaos}
referred to the absence of eventual periodicity, which bounded-error quasi-periodicity with
irrational rotation number precisely explains. For the higher families that govern the
census above, the tracks appear to be multi-branch and their arithmetic remains open; the adjacent-pair phenomenon is \emph{not} a single crossing lingering: the bite rows show
that adjacent doubles arise from different family pairs ($7\times29$ crosses
$(5,2)$--$(6,1)$ while $7\times30$ crosses $(3,4)$--$(4,3)$; $8\times22$ crosses
$(1,7)$--$(2,6)$ while $8\times23$ crosses $(6,2)$--$(7,1)$), so the clustering of
crossings at adjacent widths across \emph{distinct} track pairs is itself the phenomenon
demanding explanation. Whether triples occur infinitely often --- and where the next one is --- now looks
like a question about simultaneous approximation of these rotation numbers.

\subsection*{Acknowledgements}
I thank Shalosh B. Ekhad and Doron Zeilberger for posing the challenge and for the tables
that calibrated this work, and Doron Zeilberger and Manuel Kauers for arXiv endorsements.
All code, verification logs, and track data are available at the repository cited above.

\begin{thebibliography}{9}
\bibitem{Gale} D.~Gale, A curious Nim-type game, \emph{Amer.\ Math.\ Monthly} 81 (1974), 876--879.
\bibitem{EZ} S.~B.~Ekhad and D.~Zeilberger, All the winning bites for $a$ by $b$ Chomp for $a$ and $b$ up to 14 and two computational challenges, Personal Journal of S.~B.~Ekhad and D.~Zeilberger, Aug.~11, 2018.
\bibitem{BHMS} A.~E.~Brouwer, G.~Horv\'ath, I.~Moln\'ar-S\'aska, and C.~Szab\'o, On three-rowed Chomp, \emph{Integers} 5 (2005), \#G07.
\bibitem{Sheiner} E.~Sheiner, Unique winning opening move in three-row Chomp, arXiv:2605.23837 (2026).
\bibitem{Garg} A.~Garg, Structural results for $4\times n$ Chomp: unique extension, bimodal asymptotic structure, and period-112 geometry, arXiv:2604.25952 (2026).
\bibitem{Zchaos} D.~Zeilberger, Chomp, recurrences and chaos(?), \emph{J.\ Difference Equ.\ Appl.} 10 (2004), 1281--1293.
\bibitem{FL} E.~J.~Friedman and A.~S.~Landsberg, Nonlinear dynamics in combinatorial games: renormalizing Chomp, \emph{Chaos} 17 (2007), 023117.
\bibitem{Byrnes} S.~Byrnes, Poset game periodicity, \emph{Integers} 3 (2003), G3.
\end{thebibliography}
\end{document}
